\clearpage
\subsubsection{Rectilinear representation of phases and return}
\label{mono-recta:desarrollo}

The representation on the segment \([0,10]\) preserves the positional
origin of APP and orders its nine interior marks. In this diagram,
the integers \(1,\ldots,9\) are phase representatives. Horizontal
position indicates the order of the charts; upper labels
specify regional operations. Figure
\ref{mono-recta:figura} thus brings together the initial line, fixation of the
self-scaling axis, the specular distribution of the two channels, and
return with memory. The axis and aggregate relations and the regional
census are made explicit below; return is governed by
Proposition~\ref{mono-ampl:vueltas}.

% Recovery of fig_nonadic_route_mini.tex from the academic synthesis.
% Labels designate regional operations; the sum is performed in F_3^6.
\begin{figure}[H]
\centering
\begin{tikzpicture}[x=1.12cm,y=1cm,font=\small,>=Latex,
  guide/.style={line width=.45pt,black!55}]
 \draw[line width=.8pt] (0,0)--(10,0);
 \foreach \k in {0,...,10}{
  \draw[guide] (\k,-.09)--(\k,.09);
  \node[below=2.5mm] at (\k,0) {$\k$};
 }
 \foreach \k in {1,...,9}{\fill (\k,0) circle (1.7pt);}
 \node[anchor=west,font=\footnotesize] at (0,3.05)
   {Generating segment $[0,10]$ and nine phase representatives};
 \draw[guide] (5,.12)--(5,1.7);
 \node[above,align=center,font=\footnotesize] at (5,1.7)
   {axis fixation\\$u^\varphi$};
 \draw[decorate,decoration={brace,amplitude=4pt},guide]
   (6,.65)--(8,.65);
 \node[above=5pt,align=center,font=\footnotesize] at (7,.65)
   {regional aggregate\\$u^e+u^\pi$};
 \draw[guide] (9,.12)--(9,1.7);
 \node[above,align=center,font=\footnotesize] at (9,1.7)
   {specular involution\\$\pi\leftrightarrow e$};
 \draw[->,line width=.7pt] (9,-.7)
   .. controls (9,-1.55) and (1,-1.55) .. (1,-.7);
 \node[fill=white,inner sep=3pt,font=\footnotesize] at (5,-1.23)
   {return $9\to1$;\quad $m\mapsto m+1$};
 \node[font=\footnotesize] at (5,-1.95)
   {$w_6\to w_{12}\to w_{18}\to w_{24}\to w_{30}\to R_{36}\to G_9$};
\end{tikzpicture}
\caption{Rectilinear representation of the nonadic connection. The mark
\(5\) locates the first strong saturation and the regional
self-scaling axis; the brace \(6\)--\(8\) identifies the ternary aggregate
fixed within each fiber; phase \(9\) selects the closure
orientation. Return preserves phase and increases memory. The
abscissae number operational positions, while the labels
\(u^\pi,u^e,u^\varphi\) designate blocks of \(\mathbb F_3^6\).}
\label{mono-recta:figura}
\end{figure}


Fixation of the axis is formulated in the fiber of a given boundary. If
\(q=u^\pi+u^e+u^\varphi\) and
\(a=u^\pi+u^e-u^\varphi\), then
\[
 u^\varphi=2(q-a),\qquad s=u^\pi+u^e=q-u^\varphi
 \quad\text{in }\mathbb F_3^6.
\]
The data \((q,a)\) determine the axis and aggregate; the ordered distribution
of \(s\) retains the specular freedom resolved by prolongation.
The fifth phase has a local solution and an extension, with boundary
datum \(c=111111\) and residual signature \(\rho_W=000\). This is the
meaning of its first strong saturation. The symbol \(\varphi\)
placed above the mark \(5\) identifies the regional axis entering
that operation.

The intermediate phases preserve this decomposition within each fiber,
while their boundary data evolve. The complete census of local
multiplicities and extensions of the distinguished branch is
\[
\begin{array}{c|rrrrrrrrr}
\text{phase}&1&2&3&4&5&6&7&8&9\\\hline
N_{\rm loc}&3&2&5&4&1&1&3&3&2\\
N_{\rm ext}&2&5&4&1&1&3&3&2&1
\end{array}
\]
Phase \(6\) combines local uniqueness and three prolongations; phase
\(7\) resolves a triple specular fiber; phase \(8\) synchronizes
the three regional censuses. Phase \(9\) preserves two local
distributions with the same axis and aggregate, one of which
prolongs. These counts express different functions within a
single connection and justify the marks on the line.

In particular, for phases \(6,7,8\), the regional sums are
respectively \(012100\), \(101001\), and \(112000\). The conservation
represented by the brace is an invariance under the specular
distribution within each fiber. Subsequent real evaluation applies the
readers of compatible cylinders to complete histories; the
internal sum represented in the figure belongs to the ternary domain.

Finally, the arc \(9\to1\) represents passage from one turn to the
next. The formulas of Section
\ref{mono-ampl:conexion} preserve the phase remainder and increase the
turn counter. The rectilinear diagram and the cyclic diagram in
Figure \ref{mono-ampl:fig-regional} therefore express two complementary
aspects of the same holonomy: the order of positions and
composition of the return on the state with memory.
